\section{Final Remarks and Future Directions}
This is the first work exploring the interplay between bandit learning and online packet scheduling with deadlines. We introduce the $K$-OPSD problem, a theoretical framework that is motivated by applications such as QoS buffer management and digital advertising. We show that $K$-OPSD strictly generalizes the $K$-armed sleeping bandit problem, and develop a theory that naturally extends existing algorithms from sleeping bandits to $K$-OPSD.
%To do so, we face multiple technical challenges.

%\subsection{Challenges and Techniques}
There are several criticalities that emerge when trying to combine techniques from the toolbox of regret minimization in MABs and the analysis techniques for competitive ratio.
\paragraph{Buffer Divergence between \alg and \opt} In MABs, an algorithm can always choose the same action as \opt. This is not the case in $K$-OPSD. A single different decision leads to different buffers in the subsequent rounds. Thus, it is not possible, in general, to bound the instantaneous regret of an algorithm in this setting. Bounding the $\alpha$-regret in $K$-OPSD requires a charging scheme as an intermediate step that maps a set of subsequent actions from the algorithm to a set of actions in \opt schedule. This issue is particularly critical in randomized algorithms, where the stochastic structure exhibits an interplay between the randomness of the environment and the one of the algorithm. 
\paragraph{Breaking the $\Phi$ barrier for the competitive ratio} In the $2$-bounded OPSD setting, when the number of distinct weights is finite and equal to $K$, we provide an algorithm that achieves a competitive ratio of $\theta_K < \Phi$. The values of $\theta_K$ span from $\sqrt{2}$ to $\Phi$ as $K$ grows from $2$ to $\infty$. The algorithm builds over the well-known family of $\beta$-Earliest Deadline First ($\beta$-EDF) algorithms. $\beta$-EDF uses a static threshold to decide how close to the expiration the next packet to schedule must be. Our algorithm, \algtheta, uses a dynamic threshold and a restart mechanism. The sequence of thresholds is the solution to a system of $K+1$ equations that converge to the lower bound instance from \cite{hajek2001competitiveness}. Using a dynamic threshold requires a carefully crafted charging mechanism, where the schedule of \algtheta is divided in epochs. We write the competitive ratio as the solution to a linear optimization problem, and search for the worst possible vertex of the solution space.
\paragraph{Providing a $\theta_{K-1}$-regret lower bound} We provide a lower bound for the $\theta_{K-1}$-regret in $2$-bounded $K$-OPSD instances of $\Omega\left(\sqrt{T}\right)$. Proving a lower bound on the regret in bandits usually requires two instances that are hard to distinguish (from a statistical perspective) but with different enough optimal strategies. On the other hand, proving a lower bound on the competitive ratio for OPSD requires a carefully crafted instance where every decision always leads to the same competitive ratio. The two objectives can easily collide: two instances where every decision always leads to the same outcome cannot be different enough to generate a meaningful regret lower bound. We overcome this problem by proposing a \textit{three-instance} construction where the arrival sequence is partitioned into gadgets that repeats.
\paragraph{UCBs are not enough for randomized algorithms}
Our algorithms highlight a crucial difference in the construction of deterministic learning algorithms from randomized ones. Deterministic learning algorithms only use the notion of \textit{upper confidence bound} (UCB) to estimate the expected value of a packet type. Being optimistic is enough to upper bound the $\alpha$-regret in this case. Instead, our randomized algorithms make use of both UCBs and lower confidence bounds (LCBs). Indeed, LCBs allow randomized algorithms to deal with ratios of estimates in a more natural way. This emerges from the proofs of Theorems \ref{thr:upperboundrtwo} and \ref{thr:upperboundrs}.

% Our algorithm \algtheta improves over the best existing competitive ratio for deterministic algorithms when $K<\infty$: this result is of independent interest for the OPSD literature as it generalizes the well-known $\texttt{EDF}_{\beta}$ approach. We prove that nearly-optimal learning can be done with deterministic algorithms in $2$-bounded $K$-OPSD, providing matching upper and lower bounds. Finally, we shift our focus to randomized algorithms and propose two novel algorithms achieving $\alpha$-no-regret with a competitive ratio that matches the literature state-of-the-art.
% We study bandit learning in the Packet Scheduling with Deadlines problem. First, we show that using a reduction to sleeping bandits, it is possible to achieve sublinear $C_s$-regret for $s$-Bounded Delay $K$-OPSD, where $C_s \in [\Phi, 2)$, and $C_s = \Phi$ for $s\in \{2, 3\}$. Then, we explore two ways to beat the $\Phi$ competitive ratio: (a) using the finite $K$ types for deterministic algorithms, and (b) using randomisation. For deterministic algorithms under finite $K$ types, and $2$-Bounded Delay, we propose a novel algorithm \algtheta that beats the $\Phi$ competitive ratio without learning. This algorithm can be of independent interest to the packet scheduling community. We also propose a learning version \algthetalearn that achieves an $O(\sqrt{KT})$ $\theta_K$ regret. We also complement this section with an $\Omega(\sqrt{T})$ lower bound on the worst-case $\theta_k$-regret of any algorithm, proving the tight dependence on $T$ of our $\theta_K$-regret. For randomized algorithms, we achieve sublinear $\frac{5}{4}$-regret in $2$-Bounded instances, and sublinear  $\frac{e}{e-1}$-regret in $s$-bounded Instances for $s> 2$. All our results show that learning is possible in packet scheduling problems using variants based on upper and lower confidence bounds on the real weights.
\subsection{Future Directions}
Many interesting future directions can be explored. First, for deterministic algorithms, providing tight $\Phi$-regret for general instances is a challenging problem. Also, combining both randomized algorithms and the $K$-finite type assumptions can be explored to have competitive ratios even closer to $1$. We conjecture that the competitive ratio can be lowered below $\frac{5}{4}$, up to $\frac{7}{6} =  1.1\bar{6}$ when $K=2$, employing a similar lower-bound matching algorithm over the construction Theorem 4 of \cite{chin2003online}. Finally, additional inter-packet structure can be further explored, developing models resembling linear bandits and/or contextual bandits. All these steps towards generalization are interesting and technically challenging, but can be well-motivated by real-world applications.

% Summary: introducin $K$-OPSD learning problem. 
% Providing sublinear $C_s$-regret for general $s$-Bounded, where $C_s \in [\phi, 2)$
% Breaking the $\Phi$ barrier, utilising the finite $K$ type assumptions for deterministic algorithms, and randomisation.

% Limitations (needed by the neurips checklist):

% Future work: Adding structure to the learning problem, to go beyond the finite $K$ types assumption.






%% To add in acks 
% We thank Jean-Hugues Doussain (Criteo) for useful discussions on the motivating problem.